% Identity Bedrock — LaTeX source
% Target: IACR Cryptology ePrint Archive
% Compile: pdflatex -> bibtex -> pdflatex -> pdflatex

\documentclass[11pt,a4paper]{article}

\usepackage[utf8]{inputenc}
\usepackage[T1]{fontenc}
\usepackage{lmodern}
\usepackage{microtype}
\usepackage{amsmath,amssymb,amsthm}
\usepackage{mathtools}
\usepackage{booktabs}
\usepackage{array}
\usepackage[colorlinks=true,linkcolor=black,citecolor=black,urlcolor=blue]{hyperref}
\usepackage{xcolor}
\usepackage{listings}
\usepackage{geometry}
\usepackage{fancyhdr}
\usepackage{url}
\usepackage{cite}
\usepackage{graphicx}

\geometry{a4paper,top=25mm,bottom=25mm,left=25mm,right=25mm}

\newtheorem{theorem}{Theorem}[section]
\newtheorem{lemma}[theorem]{Lemma}
\newtheorem{proposition}[theorem]{Proposition}
\newtheorem{corollary}[theorem]{Corollary}
\theoremstyle{definition}
\newtheorem{definition}[theorem]{Definition}
\newtheorem{remark}[theorem]{Remark}

\lstset{basicstyle=\small\ttfamily,breaklines=true,frame=single,
        numbers=none,xleftmargin=1em,language={}}

\pagestyle{fancy}
\fancyhf{}
\lhead{\small Identity Bedrock}
\rhead{\small Bouchard, 2026}
\cfoot{\thepage}
\renewcommand{\headrulewidth}{0.4pt}

\newcommand{\OWF}{\mathrm{OWF}}
\newcommand{\SIG}{\mathrm{SIG}}
\newcommand{\EUFCMA}{\mathrm{EUF\text{-}CMA}}
\newcommand{\negl}{\mathrm{negl}}
\newcommand{\PBAT}{\textsf{PBAT}}
\newcommand{\BD}{\textsf{Identity Bedrock}}

\begin{document}

\title{\textbf{Identity Bedrock: Operationalizing Rompel's Minimal
Cryptographic Substrate\\for Sovereign Identity with Physical Presence Proof}}

\author{
  Jean-Fran\c{c}ois Bouchard\\[4pt]
  \small PrivateNess Network (co-founder, CTO)\\
  \small Emercoin (community contributor)\\[4pt]
  \small \texttt{[contact via privateness.network]}
}

\date{March 2026 \quad|\quad Version 1.0}

\maketitle

\begin{center}
\small\itshape
\textbf{AI disclosure:} The author used AI language model assistance
for manuscript drafting and editing. All technical contributions,
system design, implementation, theoretical framing, and binding
insights described herein originate with the author, drawing on nine
years of operational experience deploying Emercoin-based infrastructure.
\end{center}

\vspace{4pt}\noindent\rule{\textwidth}{0.4pt}\vspace{2pt}

\begin{abstract}
We present \BD{}, a deployed identity system grounded explicitly in
Rompel's proof that one-way functions ($\OWF$) are necessary and
sufficient for existentially unforgeable signatures under adaptive
chosen-message attack ($\SIG_{\EUFCMA}$)~\cite{Rompel90}. The system
separates identity cryptography (Ed25519 or ML-DSA) from directory
consensus via a chain-agnostic append-only ledger, eliminating
institutional keys beneath the individual. Deployed on Emercoin NVS
with Bitcoin AuxPoW anchoring, the attack cost on the directory layer
approaches five times EMC market capitalization~\cite{Kho21} due to
the triple consensus barrier.

We further identify bone conduction audio as a physical phenomenon
satisfying the inaccessible entropy property of Haitner
et~al.~\cite{Haitner10}: the mechanical impedance mismatch at the
air-bone interface ($T \approx 1.1 \times 10^{-3}$) renders the
authentication channel unconditionally inaccessible to replay,
independent of computational assumptions. We integrate this as a
fourth authentication factor and propose the Physically-Bound
Authentication Token (\PBAT{}), a signed structure carrying biometric
session evidence alongside cryptographic proof.

The contribution lies in the binding of existing primitives into a
system whose emergent security properties none of them exhibits alone.
\end{abstract}

\medskip
\noindent\textbf{Keywords:} one-way functions, digital signatures,
sovereign identity, bone conduction authentication, inaccessible
entropy, blockchain identity, post-quantum cryptography,
physical presence proof.

\noindent\rule{\textwidth}{0.4pt}
\tableofcontents
\newpage

\section{Introduction}

\subsection{The Problem: Institutional Capture}

Every deployed digital identity system places institutional keys
beneath the user: certificate authorities, identity providers, key
registries, platform operators. These entities can revoke, deny,
surveil, or control identity without the user's consent.

The threat model motivating this work assumes an adversary with
substantial conventional resources, potential political backing, and
full control of traditional identity infrastructure. The central
question is: \emph{can identity exist without institutional keys
beneath the individual?}

\subsection{The Theoretical Answer: Rompel 1990}

Rompel proved that $\OWF$s are both necessary and sufficient for
$\SIG_{\EUFCMA}$ secure signature schemes~\cite{Rompel90}. This is
the provable minimum --- nothing strictly weaker supports public
authentication under standard assumptions.

Our thesis: fixing $\OWF$ as the cryptographic floor and adding only
a chain-agnostic append-only directory produces an identity
construction that minimizes institutional capture to zero beneath
the individual key.

\subsection{The Binding Insight}

This work follows a pattern established by prior binding innovations
in cryptographic systems: the value lies not in the invention of new
primitives but in recognizing that existing primitives, bound in a
specific way, produce emergent properties none of them exhibits alone.
Each component of \BD{} --- $\OWF$-based signatures, append-only
directories, bone conduction biometrics --- exists independently in
prior literature. The contribution is the binding and the formal
argument for the security properties that binding produces.

\subsection{Contributions}

\paragraph{Contribution 1 --- Theory.}
Formalization of a Rompel-grounded, chain-agnostic identity system.
We prove that removing the $\OWF$ floor collapses authentication, and
that weakening the directory trades one consensus for another without
improving the institutional capture property.

\paragraph{Contribution 2 --- Systems.}
Architecture: Ed25519 or ML-DSA~\cite{FIPS204} identity layer with
pluggable directory substrate. Implementation: Emercoin NVS
(Bitcoin-anchored via AuxPoW)~\cite{Emercoin} and NESS
Fiber~\cite{DPO} (announced publisher keys, append-only BoltDB).
Deployed and publicly verifiable since 2013.

\paragraph{Contribution 3 --- Physical Presence Extension.}
Identification of bone conduction audio as a physical phenomenon
satisfying the inaccessible entropy property of Haitner
et~al.~\cite{Haitner10}, providing a formal argument for replay
resistance that prior biometric work demonstrates only
empirically~\cite{HeadSonic25,HCRAuth24}. Integration as a fourth
authentication factor and the \PBAT{} token specification.

\subsection{Paper Organization}

Section~\ref{sec:background} reviews background and related work.
Section~\ref{sec:threat} defines the threat model.
Section~\ref{sec:architecture} presents the \BD{} architecture.
Section~\ref{sec:implementation} describes the deployed implementation.
Section~\ref{sec:physical} introduces the physical presence extension.
Section~\ref{sec:evaluation} evaluates the system.
Section~\ref{sec:discussion} discusses limitations and future work.

\section{Background and Related Work}
\label{sec:background}

\subsection{Rompel's Theorem (1990)}

From a complexity-theoretic perspective, $\OWF$s, universal one-way
hash functions (UOWHFs), and $\SIG_{\EUFCMA}$ schemes are equivalent
in strength~\cite{NaorYung89}. Rompel~\cite{Rompel90} and later
Haitner et~al.~\cite{Haitner10} showed how to construct UOWHFs from
arbitrary $\OWF$s using inaccessible entropy. A modern exposition is
given in~\cite{KatzKoo05}. The equivalence is:
\[
  \OWF \;\Leftrightarrow\; \SIG_{\EUFCMA}
\]
under standard definitions~\cite{Rompel90,Goldreich01}. To our
knowledge, prior work does not explicitly operationalize this
equivalence as the design floor for a deployed identity system.

\subsection{Inaccessible Entropy}
\label{sec:inaccessible}

Haitner et~al.~\cite{Haitner10} introduce inaccessible entropy to
simplify and improve Rompel's construction. Informally, a function $F$
has inaccessible entropy if the real entropy of $F^{-1}$ ---the
entropy remaining in the input after observing the output---
significantly exceeds the accessible entropy: the maximum entropy any
polynomial-time adversary can generate among valid preimages.

\begin{definition}[Inaccessible Entropy, informal~{\cite{Haitner10}}]
\label{def:inaccessible}
Let $F: \{0,1\}^n \to \{0,1\}^m$. The \emph{inaccessible entropy}
of $F^{-1}$ is the gap
\[
  \Delta \;=\; H(X \mid F(X)) \;-\; \max_{A \in \mathrm{PPT}}
  H(A(X;\,R) \mid X)
\]
where $X$ is uniform on $\{0,1\}^n$, $R$ denotes uniform random
coins for $A$, and $A$ is constrained to output elements of
$F^{-1}(F(X))$.
\end{definition}

We observe in Section~\ref{sec:physical} that bone conduction
independently satisfies an analogous property --- but enforced by
physical rather than computational means, yielding a strictly stronger
guarantee.

\subsection{Blockchain-Based Identity Systems}

\textbf{Namecoin / EmerDNS.} Name-value storage on blockchain. Gap:
identity keys tied to blockchain wallet keys (not chain-agnostic);
miners act as an implicit certificate authority.

\textbf{W3C DIDs / Ethereum.} DID methods, smart contracts, ENS. Gap:
DAO governance keys, registry administrators, and provider
dependencies introduce institutional keys above or beneath the user.

\textbf{Sovrin / Hyperledger Indy.} Permissioned validators,
credential schemas. Gap: steward keys control the network.

\textbf{\BD{} distinction.} Pure $\OWF$ floor. Chain-agnostic
separation of identity cryptography from directory consensus. No
institutional key beneath the user by construction.

\subsection{Bitcoin AuxPoW and Merged Mining}

Auxiliary Proof-of-Work allows Emercoin to be merge-mined with
Bitcoin, inheriting a portion of Bitcoin's security for non-monetary
applications. Khovayko~\cite{Kho21} introduces the \emph{hack51}
coefficient---defined as $\text{attack\_cost} /
\text{market\_cap}$---as a numerically measurable metric for 51\%
attack resistance. Under this metric, Emercoin's triple consensus
(PoS + PoW + merged mining) produces a coefficient significantly
exceeding 1.0: attacking Emercoin requires resources approaching five
times its market capitalization, exceeding the cost of attacking
Bitcoin alone at equivalent market cap. We apply the \emph{hack51}
framework in Section~\ref{sec:eval_attack}.

\subsection{Bone Conduction Biometric Authentication}

Recent work demonstrates that bone conduction audio carries unique
per-individual characteristics. HeadSonic~\cite{HeadSonic25} proposes
a bone conduction earphone authentication system achieving 96.59\%
balanced accuracy across 60 subjects and demonstrating resilience
against replay and synthesis attacks. HCR-Auth~\cite{HCRAuth24}
demonstrates reliable authentication via head contact response.
These works establish \emph{empirically} that bone conduction
authentication resists replay attacks. Neither provides a formal
argument for why. We provide that argument in
Section~\ref{sec:impedance}.

\subsection{Skycoin Fiber Consensus}

Skycoin Fiber~\cite{Skycoin} uses announced block publisher keys
rather than economic consensus. Blocks must be signed by announced
keys; any number of hostile nodes up to $N-1$ cannot forge blocks
without those keys --- only the node holding the publisher key can
produce valid blocks. The append-only BoltDB backend makes history
rewriting detectable and maximally visible.

\subsection{Summary of Differences}

\BD{} differs from all prior work above in three respects
simultaneously. First, it is the only system explicitly grounded in
Rompel's $\OWF$--signature equivalence as a design floor. Second, the
chain-agnostic architecture decouples identity cryptography from
directory consensus: the same key works across directories, and
migration is authenticated by the key itself. Third, by identifying
bone conduction as satisfying the inaccessible entropy property, we
provide the first formal argument for why bone conduction biometrics
resist replay---prior work demonstrates this empirically without
naming the underlying principle. The \PBAT{} specification is the
first token format carrying physically-bound biometric session
evidence alongside a cryptographic identity proof in a single signed
structure.

\section{Threat Model and Requirements}
\label{sec:threat}

\subsection{Adversary Capabilities}

\textbf{Assumed:} Substantial conventional resources and political
backing; full control of traditional identity infrastructure (CAs,
IdPs, KYC registries); ability to outlaw software, seize property,
and coerce individuals; classical (polynomial-time) computational
adversary.

\textbf{Not assumed:} Cannot break $\OWF$ under standard complexity
assumptions; cannot mount a sustained 51\% attack on Bitcoin
(quantified in Section~\ref{sec:eval_attack}); cannot simultaneously
compromise distributed infrastructure globally; cannot violate
acoustic physics (Section~\ref{sec:impedance}).

\textbf{Explicit scope.} This work addresses classical adversaries.
Post-quantum migration to ML-DSA (FIPS~204)~\cite{FIPS204} is
discussed in Section~\ref{sec:limitations}.

\subsection{Design Requirements}

\textbf{R1 --- Mathematical irreducibility.} $\OWF$ as cryptographic
floor. No strictly weaker primitive beneath.

\textbf{R2 --- No institutional keys.} Individual key is sole
cryptographic root of trust. No CA, IdP, or registry key beneath
the user.

\textbf{R3 --- Chain-agnostic.} Identity cryptography independent of
any specific directory consensus. Portable across append-only
substrates.

\textbf{R4 --- Quantifiable security.} Attack costs explicit and
measurable.

\textbf{R5 --- Operational verification.} Deployed system, not
theoretical proposal. Public verification of all claims.

\textbf{R6 --- Physical presence.} For deployments requiring physical
presence proof, the system must provide it unconditionally --- not
via algorithmic liveness detection but via physical channel
properties.

\section{Identity Bedrock Architecture}
\label{sec:architecture}

\subsection{The Two-Layer Design}

\paragraph{Layer 1: Identity Cryptography (Chain-Agnostic).}
\[
\text{Seed} \;(\geq 256\text{ bits}) \;\xrightarrow{\mathsf{KeyGen}}\;
(\mathit{sk},\; \mathit{pk})
\]
\[
\text{Identity} = (\mathit{id},\, \mathit{pk}), \qquad
\text{Authentication} = \mathsf{Sign}_{\mathit{sk}}(\mathit{challenge})
\]

Properties: $\OWF$-equivalent security (Rompel equivalence);
independent of any blockchain; portable across directory substrates.
Default instantiation: Ed25519~\cite{RFC8032}. Migration target:
ML-DSA-65~\cite{FIPS204}.

\paragraph{Layer 2: Directory Substrate (Pluggable).}

\begin{definition}[Append-Only Directory with Finality]
\label{def:directory}
A directory $\mathcal{D} \subseteq \mathcal{ID} \times \mathcal{PK}$
is \emph{append-only with finality} if: \emph{(i)} once
$\mathsf{Bind}(\mathit{id}, \mathit{pk}) = 1$ is confirmed, it cannot
be set to $0$ without an attack of quantifiable cost; and \emph{(ii)}
the binding is globally verifiable by any party without trusting an
intermediary.
\end{definition}

\paragraph{Separation invariant.} Identity uses its chosen signature
scheme; the directory uses whatever consensus fits the deployment's
threat model. The same identity key can move between directories.
No coupling exists between the two layers.

\subsection{Authentication Protocol}

\paragraph{Setup.}
\begin{enumerate}
  \item User generates $(\mathit{sk}, \mathit{pk})$ from entropy source.
  \item User publishes $(\mathit{id}, \mathit{pk})$ to directory
        $\mathcal{D}$.
  \item Wait for finality (directory-dependent).
  \item Binding is global: $\mathsf{Bind}(\mathit{id}, \mathit{pk})=1$.
\end{enumerate}

\paragraph{Challenge-Response.}
\begin{enumerate}
  \item $V \to P$: fresh nonce $c \xleftarrow{\$} \{0,1\}^\lambda$.
  \item $P \to V$: $\sigma = \mathsf{Sign}_{\mathit{sk}}(c)$.
  \item $V$ verifies:
    \begin{enumerate}
      \item Fetch $\mathit{pk}$ s.t.\
            $\mathsf{Bind}(\mathit{id}, \mathit{pk}) = 1$ from
            $\mathcal{D}$.
      \item Check $\mathsf{Verify}_{\mathit{pk}}(c, \sigma) = 1$.
      \item Accept iff both hold.
    \end{enumerate}
\end{enumerate}

No state beyond $(\mathit{id}, \mathit{pk})$ is held in the directory.
No smart contracts. Verification is pure mathematics.

\subsection{Security Proof}

\begin{theorem}
\label{thm:unforgeability}
If the underlying signature scheme is $\SIG_{\EUFCMA}$ secure and
$\mathcal{D}$ is append-only with finality
(Definition~\ref{def:directory}), then \BD{} provides unforgeable
authentication against any classical adversary.
\end{theorem}

\begin{proof}
Suppose adversary $\mathcal{A}$ wins the authentication game: it
produces valid $(\mathit{id}^*, c^*, \sigma^*)$ without holding
$\mathit{sk}$.

Construct forger $\mathcal{B}$ against the signature scheme.
$\mathcal{B}$ receives challenge public key $\mathit{pk}^*$. It
simulates the directory by answering binding queries for
$\mathit{id}^*$ with $\mathit{pk}^*$.\footnote{$\mathcal{B}$
simulates the directory locally; it does not write to any real
ledger.} When $\mathcal{A}$ triggers authentication, $\mathcal{B}$
forwards the $\SIG_{\EUFCMA}$ game challenge $c^*$ to $\mathcal{A}$.
Since $c^*$ was chosen fresh and never signed by the legitimate
holder of $\mathit{sk}$, $\mathcal{A}$'s forgery $\sigma^*$
constitutes an existential forgery. By security of the scheme:
\[
  \Pr[\mathcal{A}\text{ wins}] \;\leq\; \negl(\lambda).
\]
Directory attacks shift cost to the directory layer, quantified in
Section~\ref{sec:eval_attack}.\hfill\qed
\end{proof}

\subsection{Chain-Agnostic Property and Migration}

The same $(\mathit{sk}, \mathit{pk})$ can be published to Emercoin
NVS, Bitcoin OP\_RETURN, Ethereum contract storage, or any
append-only ledger satisfying Definition~\ref{def:directory}.
Migration is authenticated: the user signs a migration statement with
$\mathit{sk}$, publishes to $\mathcal{D}'$, and verifiers accept
after finality. The old directory retains the historical record.

\section{Implementation: Emercoin + NESS Fiber}
\label{sec:implementation}

The architecture described in this paper builds on the PrivateNess
Network design published in October 2020~\cite{Litepaper20}, which
established the dual-blind Emercoin + Skywire architecture, 127.0.0.1
communication restriction, UHE-PRNG entropy at 1536 bits of internal
state, and the full OSI-stack decentralization approach. Identity
Bedrock formalizes the cryptographic foundations of that deployed
system.

\subsection{Directory: Emercoin NVS}

Emercoin NVS provides a name-value store on a hybrid PoS + AuxPoW
blockchain~\cite{Emercoin}. Consensus: 85\% Proof-of-Stake + 15\%
Bitcoin AuxPoW. Finality: 1 confirmation ($\approx$2 minutes).

Identity bindings are anchored using the WORM (World Object MaRkup
language) schema~\cite{WORM}, a decentralized XML-based search
infrastructure designed by Aleksej Sokolov and deployed within the
PrivateNess ecosystem. The \texttt{<user>} element is the canonical
container:

\begin{lstlisting}[language=XML]
<worm>
  <user type="ness" key="<pubkey_hex>" nonce="<nonce>">
    <keys>
      <key hash="<pubkey_hash>" algo="ed25519">
        <pubkey_hex>
      </key>
    </keys>
    <node type="ness" url="<skywire_pubkey>"
          nonce="<nonce>" tags="recorder,voice"/>
  </user>
</worm>
\end{lstlisting}

The \texttt{<keys>} list natively supports multiple algorithms,
enabling ML-DSA migration without schema changes.

\paragraph{Separation of concerns.} The Emercoin wallet key controls
NVS record updates. The identity key authenticates the identity
itself. Directory control is strictly separated from identity
authentication.

\subsection{Infrastructure: NESS Fiber}

NESS Fiber uses Skycoin-type Fiber consensus~\cite{Skycoin} with
announced publisher keys published in Emercoin NVS under
\texttt{dpo/PrivateNESS.Network}~\cite{DPO}, creating a verifiable
trust chain: NESS Fiber security anchors to Emercoin NVS, which
anchors to Bitcoin via AuxPoW. This anchoring is \emph{trustless}:
any verifier can confirm the binding independently at the Emercoin
block explorer without relying on any intermediary.

Key compromise is maximally detectable: a stolen publisher key cannot rewrite past blocks (append-only BoltDB makes history rewriting impossible). The only action available is a reset to genesis block~0, which is immediately visible to all participants. Key rotation is itself a publicly auditable event.

\subsection{Canonical Self-Spend Proof}

The canonical self-spend proof binds a chain address to the Ed25519
identity:
\begin{enumerate}
  \item Create unsigned self-spend transaction $A \to A$
  \item Sign with the secp256k1 chain key
  \item Create proof $\{$\texttt{network, address, policy, txid,
        tx\_hex}$\}$
  \item Anchor proof in Emercoin NVS WORM record
  \item WORM record controlled by the Ed25519 bedrock key
\end{enumerate}

\paragraph{Deployed instance.}
The NESS Fiber chain instantiates this construction at genesis.
The genesis address
(\texttt{24GJTLPMoz61sV4J4qg1n14x5qqDwXqyJJy}),
blockchain public key
(\texttt{02933015bd2fa1...ccaf525036}),
and genesis block signature are published in
\texttt{dpo/PrivateNESS.Network}~\cite{DPO}, anchored in Emercoin
block~\#445669 (November~1, 2020, 22:16:45~UTC) via AuxPoW:
\[
\text{Bitcoin AuxPoW}
\;\to\; \text{Emercoin block \#445669}
\;\to\; \text{DPO record}
\;\to\; \text{BlockchainPubkeyStr}
\;\to\; \text{GenesisSignatureStr}
\;\to\; \text{NESS Fiber chain}
\]
Any verifier can confirm at
\url{https://explorer.emercoin.com/dpo/PrivateNESS.Network} and
\url{https://explorer.emercoin.com/block/445669}.

\subsection{Key Generation and Entropy Provisioning}

Identity keys are derived from a seed produced by Gibson's UHEPRNG~%
\cite{Gibson}. Entropy sources mixed in include: multiple 1024-byte
blocks from \texttt{os.urandom} ($\geq 1536$ bits total); RC4OK
output from Emercoin core RNG; CPU timing jitter; and hardware
identifiers (eMMC CID, MAC address, machine-id). The Linux kernel
caps its DRBG state to approximately 256 bits; our multi-source
mixing increases robustness against partial bias without claiming to
exceed this ceiling. Security rests on the standard
$\approx 128$-bit classical security floor of the chosen scheme.

\section{Physical Presence Extension}
\label{sec:physical}

\subsection{Motivation}

The two-layer \BD{} system addresses forgery ($\OWF$ layer) and
directory rewrite (AuxPoW layer). It does not address physical
presence: a device with a stolen key can impersonate the legitimate
owner. For deployments where physical presence is a security
requirement, an additional unconditionally guaranteed factor is
needed.

\subsection{The Channel Symmetry Problem}

Every voice authentication system built on air-conducted audio shares
a fundamental vulnerability: the authentication and replay channels
are physically identical. A microphone captures air pressure waves.
A speaker produces air pressure waves. This is not a software
deficiency --- it is a property of air-conducted sound that no
liveness detection algorithm can fully overcome.

\subsection{Bone Conduction Breaks Channel Symmetry}
\label{sec:impedance}

Bone conduction transmits mechanical vibration through skull bone and
soft tissue to the transducer. Replay requires producing the same
mechanical vibration pattern at the transducer --- requiring the same
skull geometry, bone density, tissue distribution, and coupling
characteristics as the enrolled individual.

The physical barrier to replay is the mechanical impedance mismatch
at the air-bone interface. The specific acoustic impedance of air is
$Z_\text{air} \approx 415$~Pa$\cdot$s/m; cortical bone is
$Z_\text{bone} \approx 1.5 \times 10^6$~Pa$\cdot$s/m~\cite{Acoustics}.
The pressure transmission coefficient is:

\begin{equation}
\label{eq:transmission}
T = \frac{4\, Z_\text{air}\, Z_\text{bone}}{(Z_\text{air} +
Z_\text{bone})^2}
= \frac{4 \times 415 \times 1.5 \times 10^6}{(415 + 1.5 \times
10^6)^2}
\approx 1.1 \times 10^{-3}
\end{equation}

Approximately 0.11\% of acoustic energy from an air-conducted source
transmits into bone. The spatial coupling characteristics --- which
depend on the geometry of the skull in contact with the transducer
--- are entirely absent in any replay attempt.

\begin{remark}
Equation~(\ref{eq:transmission}) is a consequence of acoustic
physics, not a hardness assumption. It holds unconditionally,
independent of any computational model.
\end{remark}

\subsection{Bone Conduction as Physical Inaccessible Entropy}

\begin{remark}[Physical instantiation of inaccessible entropy]
\label{rem:physical_ie}
The bone conduction channel produces a distribution --- the mechanical
resonance signature of the enrolled skull --- whose entropy is
inaccessible to any adversary not in physical control of the enrolled
body in contact with the enrolled transducer. The inaccessibility is
enforced by Equation~(\ref{eq:transmission}), not by any
computational hardness assumption. Where Definition~\ref{def:inaccessible}
requires inaccessibility against PPT adversaries (a computational
bound), bone conduction provides inaccessibility against \emph{all}
adversaries lacking physical control (an unconditional bound) ---
a strictly stronger guarantee under any computational model.
\end{remark}

\begin{remark}[Scope of Remark~\ref{rem:physical_ie}]
Remark~\ref{rem:physical_ie} is an observation connecting a physical
phenomenon to the inaccessible entropy framework of Haitner
et~al.~\cite{Haitner10}. A formal reduction demonstrating that bone
conduction channel entropy satisfies Definition~\ref{def:inaccessible}
under the standard model remains for future work.
\end{remark}

\subsection{The Four-Factor Authentication Stack}

\begin{table}[h]
\centering
\caption{Four-factor authentication stack}
\label{tab:fourfactor}
\begin{tabular}{llll}
\toprule
Factor & Mechanism & Attack addressed & Guarantee \\
\midrule
Ed25519 / $\OWF$ & Rompel floor & Forgery & Computational \\
EmerNVS / AuxPoW & Bitcoin-anchored dir. & Dir.\ rewrite & Economic \\
Bone conduction & Physical coupling & Replay & Unconditional \\
Voiceprint + \PBAT{} & Session embed.\ signed & Impersonation &
Biometric + crypto \\
\bottomrule
\end{tabular}
\end{table}

\subsection{The Physically-Bound Authentication Token}

\begin{definition}[\PBAT{}]
\label{def:pbat}
A \emph{Physically-Bound Authentication Token} is a signed data
structure containing: an identity public key; a directory binding
proof; a biometric session embedding captured via a physically
inaccessible channel (Definition~\ref{def:physical_channel}); a
single-use nonce; a timestamp; and a digital signature over all
fields by the identity private key.
\end{definition}

\begin{lstlisting}
{
  "identity_pubkey":   "<Ed25519 or ML-DSA pubkey hex>",
  "identity_algo":     "ed25519 | mldsa65",
  "nvs_record":        "<EmerNVS WORM txid>",
  "voiceprint_hash":   "<SHA256 or SHA3-256 of enrolled embedding>",
  "session_embedding": "<current session voice embedding>",
  "similarity_score":  0.94,
  "capture_channel":   "bone_conduction_msbc",
  "headset_mac":       "<BT MAC of enrolled headset>",
  "timestamp":         "<UTC ISO8601>",
  "nonce":             "<random 256-bit>",
  "expires_seconds":   300,
  "signature":         "<Sign_sk(all above fields)>"
}
\end{lstlisting}

A \PBAT{} cannot be replayed: the nonce is single-use, the timestamp
is signed, and the session embedding is captured live at this
authentication event.

\paragraph{Physical presence doctrine.}
\emph{The device authenticates the body, not the credential. No
network path, no software override, and no root privilege substitutes
for physical presence of the enrolled body wearing the enrolled
transducer.}

\paragraph{Quantum note.}
The physical presence layer is unconditionally quantum-resistant.
Post-quantum migration is scoped to the cryptographic binding layer
only.

\section{Evaluation}
\label{sec:evaluation}

\subsection{Attack Cost Analysis}
\label{sec:eval_attack}

\begin{table}[h]
\centering
\caption{Attack cost analysis}
\label{tab:attacks}
\renewcommand{\arraystretch}{1.2}
\begin{tabular}{p{2.8cm}p{3.5cm}p{3.2cm}p{3cm}}
\toprule
Attack vector & Requirements & Estimated cost & Result \\
\midrule
Break signature scheme
  & Violate $\OWF$ assumption
  & Infeasible (classical)
  & Authentication collapses \\
Rewrite Emercoin NVS
  & PoS majority + PoW + AuxPoW 51\% simultaneously
  & $\approx\!5\times$ EMC mkt.\ cap~\cite{Kho21};
    Bitcoin component $\sim$\$5--20B~\cite{CoinMetrics24}
  & Directory rewrite \\
Compromise NESS Fiber
  & Steal publisher keys
  & High; DoS only without keys
  & Cannot alter identity bindings \\
Replay bone conduction
  & Reproduce mechanical skull signature via air
  & Physically infeasible ($T \approx 10^{-3}$,
    Eq.~\ref{eq:transmission})
  & Biometric layer holds \\
Coerce enrolled user
  & Physical
  & Variable
  & Outside cryptographic scope \\
\bottomrule
\end{tabular}
\end{table}

\paragraph{Note on AuxPoW attack cost.}
Khovayko~\cite{Kho21} demonstrates via the \emph{hack51} coefficient
that Emercoin's triple consensus barrier requires resources approaching
$5\times$ EMC market capitalization. The \$5--20B
figure~\cite{CoinMetrics24} represents the Bitcoin AuxPoW component;
it is a \emph{lower bound} on the total attack cost.

\subsection{Comparison to Related Systems}

\begin{table}[h]
\centering
\caption{Comparison of identity systems}
\label{tab:comparison}
\begin{tabular}{lcccc}
\toprule
System & Identity floor & No inst.\ keys & Chain-agnostic
  & Physical presence \\
\midrule
\textbf{\BD{}} & Ed25519/Rompel & \checkmark & \checkmark
  & Physical (bone) \\
X.509 PKI & RSA/ECDSA & $\times$ & $\times$ & None \\
Namecoin & ECDSA & $\times$ & $\times$ & None \\
Ethereum DIDs & Varies & $\times$ & $\times$ & None \\
Sovrin & Ed25519 & $\times$ & $\times$ & None \\
FIDO2 & ECDSA & Partial & $\times$ & HW token \\
HeadSonic~\cite{HeadSonic25} & None & N/A & N/A
  & Physical (bone) \\
\bottomrule
\end{tabular}
\end{table}

\subsection{Performance}

\begin{itemize}
  \item \textbf{Emercoin NVS finality:} $\approx$2 minutes
  \item \textbf{Challenge-response:} milliseconds (no on-chain
        transaction required at verification time)
  \item \textbf{Voiceprint verification} (resemblyzer%
        \footnote{\url{https://github.com/resemble-ai/resemblyzer}},
        RK3528 class hardware): ${<}1$ second
  \item \textbf{\PBAT{} generation and signing:} ${<}10$ ms
\end{itemize}

\section{Discussion}
\label{sec:discussion}

\subsection{Limitations}
\label{sec:limitations}

\paragraph{Out of scope.}
Key compromise via malware or physical theft affects any cryptographic
system. Social engineering, phishing, and UI-layer attacks are
orthogonal to the security model.

\paragraph{Quantum adversary.}
Ed25519 is vulnerable to Shor's algorithm. This work scopes to
classical adversaries. The intended migration path is ML-DSA-65
(CRYSTALS-Dilithium3, FIPS~204~\cite{FIPS204}):
\begin{itemize}
  \item Replace Ed25519 with ML-DSA-65 (pubkey 1952 bytes vs.\ 32;
        accommodated by EmerNVS value field)
  \item Replace SHA-256 voiceprint hash with SHA3-256 in \PBAT{}
  \item WORM \texttt{<keys>} schema natively supports multiple
        algorithms; migration adds \texttt{<key algo="mldsa65">}
        without modifying existing records
  \item Physical presence layer, \PBAT{} structure, and directory
        anchoring unchanged
  \item OpenSSL~3.5 (April~2025) includes native ML-DSA support
        in mainline Armbian Linux
\end{itemize}

\paragraph{Formal reduction for bone conduction.}
The connection between bone conduction and
Definition~\ref{def:inaccessible} is an observation
(Remark~\ref{rem:physical_ie}), not a completed theorem. A formal
reduction under the standard model remains for future work.

\paragraph{Attack cost precision.}
Absolute figures depend on current EMC market conditions and should
be treated as order-of-magnitude illustrations.

\paragraph{Coercion.}
The residual attack surface of \BD{} is coercion, enrollment
compromise, and post-quantum cryptanalysis --- each explicitly
bounded.

\subsection{On the Nature of the Contribution}

This work follows the pattern of binding innovations in applied
cryptography: taking existing primitives and identifying a binding
that produces emergent properties none of them exhibits alone. The
emergent property is an identity system with no institutional keys
beneath the individual, unconditional physical presence proof, and a
formal cryptographic floor, deployed and publicly verifiable since
2013~\cite{Emercoin} / 2020~\cite{DPO}.

\subsection{Future Work}

\begin{itemize}
  \item Formal reduction: prove bone conduction channel entropy
        satisfies Definition~\ref{def:inaccessible} under the
        standard model
  \item Full \PBAT{} security proof under the standard model
  \item Threshold signatures ($k$-of-$n$ ML-DSA) for social
        recovery without institutional escrow
  \item Demonstrate chain-agnostic portability across Bitcoin
        OP\_RETURN and Ethereum
  \item User studies on bone conduction enrollment UX and false
        rejection rate
\end{itemize}

\section*{Conclusion}
\addcontentsline{toc}{section}{Conclusion}

We have presented \BD{}, a deployed identity system grounded in
Rompel's $\OWF$--signature equivalence. The system achieves no
institutional keys beneath the individual through a chain-agnostic
architecture separating identity cryptography from directory
consensus.

We have further identified bone conduction audio as a physical
phenomenon satisfying the inaccessible entropy property of Haitner
et~al.~\cite{Haitner10} (Remark~\ref{rem:physical_ie}), providing
the physical argument for replay resistance that prior biometric work
demonstrates empirically. Integrated as a fourth authentication
factor, this produces the four-factor system of
Table~\ref{tab:fourfactor}.

The canonical on-chain identity record at~\cite{DPO} constitutes
operational proof: the genesis address, blockchain public key, and
genesis signature are anchored in Emercoin NVS via AuxPoW,
transitively anchored to Bitcoin, and independently auditable by
any party without trusting an intermediary. This is not a paper
describing what could be built. It is a paper describing what has
been running since 2020.

\section*{Acknowledgments}
\addcontentsline{toc}{section}{Acknowledgments}

The author thanks Oleg Khovayko for foundational Emercoin
infrastructure work and the \emph{hack51} attack cost framework;
Aleksej Sokolov (Alex014) for designing the WORM format and
implementing the PrivateNess toolchain (\texttt{privatenesstools},
\texttt{pyuheprng}) from the author's specifications; Brandon
``Synth'' Smietana and Moses Narrow (\texttt{moe\_narrow} /
\texttt{0pcom}), Skywire Linux packaging lead, for the Skywire mesh
transport and Fiber consensus infrastructure on which NESS is built;
and the PrivateNess Network community, whose collective work since
the project's public founding on Bitcointalk on September~1,
2020~\cite{BTCTalk20} has made this deployed system possible. The
author's own engagement with Emercoin infrastructure spans nine
years and provides the operational foundation on which \BD{} is
built.

\bibliographystyle{plain}
\begin{thebibliography}{99}

\bibitem{Rompel90}
J.~Rompel.
\newblock One-way functions are necessary and sufficient for secure
  signatures.
\newblock In {\em Proc.\ 22nd Annual ACM Symposium on Theory of
  Computing (STOC)}, pages 387--394. ACM, 1990.

\bibitem{Haitner10}
I.~Haitner, T.~Holenstein, O.~Reingold, S.~Vadhan, and H.~Wee.
\newblock Universal one-way hash functions via inaccessible entropy.
\newblock In H.~Gilbert, editor, {\em Advances in Cryptology ---
  EUROCRYPT~2010}, LNCS vol.~6110, pages 616--637. Springer, 2010.

\bibitem{RFC8032}
S.~Josefsson and I.~Liusvaara.
\newblock Edwards-curve digital signature algorithm ({EdDSA}).
\newblock RFC 8032, IETF, January 2017.
\newblock \url{https://www.rfc-editor.org/rfc/rfc8032}

\bibitem{Gibson}
S.~Gibson.
\newblock Ultra-high entropy pseudo-random number generator ({UHEPRNG}).
\newblock Gibson Research Corporation.
\newblock \url{https://www.grc.com/otg/uheprng.htm}

\bibitem{Emercoin}
Emercoin Development Team.
\newblock Emercoin: hybrid blockchain with merged mining.
\newblock \url{https://emercoin.com}, 2013--2026.

\bibitem{Kho21}
O.~Khovayko.
\newblock {Оценка стоимости атаки 51\% на {Emercoin}:
  коэффициент hack51}.
\newblock Dzen / Yandex, 2021.
\newblock \url{https://dzen.ru/a/YUotZkN9nAeElPEQ} [In Russian]

\bibitem{CoinMetrics24}
L.~Nuzzi et~al.
\newblock Breaking {BFT}: quantifying the cost to attack {Bitcoin}
  and {Ethereum}.
\newblock Coin Metrics, 2024.

\bibitem{Skycoin}
B.~Smietana (``Synth'') et~al.
\newblock Skycoin: {Skywire} and {Fiber} consensus.
\newblock Skycoin Development Team, 2014--2026.
\newblock \url{https://www.skycoin.com}
\newblock Linux packaging: M.~Narrow (\texttt{moe\_narrow} /
  \texttt{0pcom}).

\bibitem{WORM}
A.~Sokolov (Alex014) / {NESS-Network}.
\newblock {WORM} --- {W}orld {O}bject {Ma}R\-kup language.
\newblock Format designed by A.~Sokolov; deployed within the
  {PrivateNess} ecosystem.
\newblock \url{https://github.com/ness-network/worm}

\bibitem{NaorYung89}
M.~Naor and M.~Yung.
\newblock Universal one-way hash functions and their cryptographic
  applications.
\newblock In {\em Proc.\ 21st Annual ACM Symposium on Theory of
  Computing (STOC)}, pages 33--43. ACM, 1989.

\bibitem{KatzKoo05}
J.~Katz and C.-Y. Koo.
\newblock On constructing universal one-way hash functions from
  arbitrary one-way functions.
\newblock Cryptology ePrint Archive, Report 2005/328, 2005.
\newblock \url{https://eprint.iacr.org/2005/328}

\bibitem{FIPS204}
{National Institute of Standards and Technology}.
\newblock Module-lattice-based digital signature standard ({ML-DSA}).
\newblock {FIPS}~204, August 2024.
\newblock \url{https://doi.org/10.6028/NIST.FIPS.204}

\bibitem{HeadSonic25}
Z.~He, J.~Chen, K.~He, Y.~Gu, Q.~Deng, and Z.~Zhang.
\newblock {HeadSonic}: usable bone conduction earphone authentication
  via head-conducted sounds.
\newblock {\em IEEE Transactions on Mobile Computing},
  24(9):7914--7928, September 2025.
\newblock doi:10.1109/TMC.2025.3551272

\bibitem{HCRAuth24}
Z.~He, J.~Chen, C.~Wu, K.~He, R.~Du, J.~Jia, Y.~Gu, and X.~Sun.
\newblock {HCR-Auth}: reliable bone conduction earphone
  authentication with head contact response.
\newblock {\em Proc.\ ACM Interact.\ Mob.\ Wearable Ubiquitous
  Technol.}, 8(4):1--27, November 2024.
\newblock doi:10.1145/3699780

\bibitem{Acoustics}
L.~E.~Kinsler, A.~R.~Frey, A.~B.~Coppens, and J.~V.~Sanders.
\newblock {\em Fundamentals of Acoustics}, 4th~ed.
\newblock John Wiley \& Sons, Hoboken, NJ, 2000.
\newblock ISBN: 978-0-471-84789-2.

\bibitem{Goldreich01}
O.~Goldreich.
\newblock {\em Foundations of Cryptography, Vol.~1: Basic Tools}.
\newblock Cambridge University Press, 2001.

\bibitem{BTCTalk20}
{PrivateNESS Network}.
\newblock {PrivateNess Network (NESS) Privacy CX Fiber Coin}.
\newblock Bitcointalk Announcements (Altcoins), September~1, 2020.
\newblock \url{https://bitcointalk.org/index.php?topic=5272741.0}

\bibitem{Litepaper20}
J.-F.~Bouchard et~al.
\newblock {PrivateNess Network LitePaper}.
\newblock PrivateNess Network, October~9, 2020.
\newblock \url{https://privateness.network/?p=255}
\newblock \textit{``The only effective weapon against big tech is
  small tech decentralized.''}

\bibitem{DPO}
{PrivateNESS Network}.
\newblock {PrivateNESS.Network} --- canonical on-chain identity
  record.
\newblock Emercoin NVS DPO, anchored in block~\#445669,
  November~1, 2020, 22:16:45~UTC. Operational 2020--2026.
\newblock \url{https://explorer.emercoin.com/dpo/PrivateNESS.Network}
\newblock Block: \url{https://explorer.emercoin.com/block/445669}

\end{thebibliography}

\appendix

\section{Formal Definitions}
\label{app:defs}

\begin{definition}[One-Way Function]
A function $f: \{0,1\}^* \to \{0,1\}^*$ is a \emph{one-way
function} if it is computable in polynomial time and for all PPT
adversaries $\mathcal{A}$:
\[
  \Pr\bigl[\mathcal{A}(f(x)) \in f^{-1}(f(x))\bigr]
  \;\leq\; \negl(\lambda)
\]
where $x \xleftarrow{\$} \{0,1\}^\lambda$.
\end{definition}

\begin{definition}[$\SIG_{\EUFCMA}$ Security]
A signature scheme $(\mathsf{Gen}, \mathsf{Sign}, \mathsf{Verify})$
is \emph{existentially unforgeable under adaptive chosen-message
attack} if for all PPT adversaries $\mathcal{A}$ with access to a
signing oracle:
\[
  \Pr\bigl[\mathsf{Verify}(\mathit{pk}, m^*, \sigma^*) = 1
  \;\wedge\; m^* \notin Q\bigr]
  \;\leq\; \negl(\lambda)
\]
where $(\mathit{pk}, \mathit{sk}) \gets \mathsf{Gen}(1^\lambda)$
and $Q$ is the set of queried messages.
\end{definition}

\begin{definition}[Physically Inaccessible Biometric Channel]
\label{def:physical_channel}
A biometric capture channel is \emph{physically inaccessible} if
the entropy produced by the channel is inaccessible to any adversary
without physical control of the enrolled subject in contact with the
enrolled transducer, by physical properties of the channel rather
than by computational hardness assumptions.
\end{definition}

\section{WORM Identity Record Schema}
\label{app:worm}

The WORM \texttt{<keys>} list supports multiple algorithms,
enabling post-quantum migration without schema changes:

\begin{lstlisting}[language=XML]
<worm>
  <user type="ness" key="<primary_pubkey_hex>"
        nonce="<random_nonce>">
    <keys>
      <key hash="<pubkey_hash>" algo="ed25519">
        <pubkey_hex>
      </key>
      <key hash="<mldsa_hash>" algo="mldsa65">
        <mldsa_pubkey_hex>
      </key>
    </keys>
    <node type="ness" url="<skywire_pubkey>"
          nonce="<skywire_nonce>"
          tags="recorder,voice"/>
  </user>
</worm>
\end{lstlisting}

\section{Deployment Verification}
\label{app:deployment}

\begin{itemize}
  \item \textbf{PrivateNess Network on-chain identity (EmerNVS DPO):}\\
        \url{https://explorer.emercoin.com/dpo/PrivateNESS.Network}\\
        Anchored in block~\#445669, November~1, 2020, 22:16:45~UTC.
  \item \textbf{Bitcointalk (September~1, 2020):}\\
        \url{https://bitcointalk.org/index.php?topic=5272741.0}
  \item \textbf{Emercoin NVS explorer:}
        \url{https://explorer.emercoin.com}
  \item \textbf{PrivateNess Network:}
        \url{https://privateness.network}
  \item \textbf{PrivateNess tools}
        (A.~Sokolov, from J.-F.~Bouchard specifications):\\
        \url{https://github.com/ness-network/privatenesstools}
  \item \textbf{pyuheprng}
        (A.~Sokolov, from J.-F.~Bouchard specifications):\\
        \url{https://github.com/ness-network/pyuheprng}
  \item \textbf{WORM} (format designed by A.~Sokolov):\\
        \url{https://github.com/ness-network/worm}
  \item \textbf{Block explorers:}\\
        \url{https://ness-explorer.magnetosphere.net}\\
        \url{https://explorer.privateness.network}
\end{itemize}

\end{document}